\DeclareMockProblem{MOCK-A-Q01}{
设 \(X\) 为紧致拓扑空间，\(A\subseteq X\) 为闭集。直接从开覆盖定义证明 \(A\) 紧致。
}{
任取 \(A\) 的开覆盖 \(\{U_i\}_{i\in I}\)，其中每个 \(U_i\) 是 \(A\) 的相对开集。
对每个 \(i\)，取 \(X\) 中开集 \(V_i\)，使 \(U_i=A\cap V_i\)。因为 \(A\) 闭，
\(X\setminus A\) 开，且
\[
  X=(X\setminus A)\cup\bigcup_{i\in I}V_i.
\]
这是一组 \(X\) 的开覆盖。由 \(X\) 紧致，可从中取有限子覆盖。删去其中可能出现的
\(X\setminus A\) 后，相应的有限多个 \(V_i\) 仍覆盖 \(A\)；于是对应的有限多个
\(U_i=A\cap V_i\) 覆盖 \(A\)。故 \(A\) 紧致。
}{
目标仍是开覆盖的有限化，但覆盖只覆盖 \(A\)，不能直接使用 \(X\) 的紧致性。要让
\(X\) 的紧致性接手，就必须把覆盖补成整个 \(X\) 的开覆盖。题目特意给出“\(A\) 闭”，
正是为了保证补集 \(X\setminus A\) 开；把这个补集加入覆盖后，\(A\) 外的所有点一次性
被覆盖，而 \(A\) 内仍由原来的开集负责。

取出有限子覆盖后，补集对覆盖 \(A\) 没有贡献，可以删去。这个构造体现了常见动机：
当已知定理作用在更大的对象上时，用一个开集处理目标外部，把局部问题提升到整个空间。
相关定义见 \knowledgelink{sec:compactness}{紧致性与连续性}。
}

